\documentclass{article}
\usepackage[margin=2cm]{geometry}
\usepackage{amsmath}
\usepackage{graphicx}
\usepackage{booktabs}
\usepackage[utf8]{inputenc}
\usepackage{ragged2e}
\setlength{\emergencystretch}{3em}
\begin{document}
\sloppy

\noindent\hfill\includegraphics[height=1.8em]{assets/figure_header_logo.png}\par\noindent\hrule\vspace{4pt}

for $C . { - \{ \mathbf { s } - \mathbf { n } \} } ^ { a }$ an arbitrary component of the master field
$C ;$ Taking ${ \mathfrak { m } } \rightarrow - { \mathfrak { m } }$
and ${ \mathfrak { s } } ~ = ~ m + 1$ in (), we can iterativelydetermine the
dependence of the $\bar { C } . { \mathfrak { m } } ^ { m + 1 }$ on the $C . { \mathbf { v } } ^ { \parallel }$.From the equation
$( )$ we obtain


The term in () that is of further interest is the one multiplying $4 e ^ { - 2 p } \partial \partial$ acting on $\delta C . { 0 } ^ { 1 }$ which is convenient to.
compute in the metric formulation..


\begin{equation}
\begin{array} { c c } { \partial _ { z } C _ { - m } ^ { m + 1 } + \frac { \mathrm { e } ^ { \mu } } { 2 } g _ { 3 } ^ { 2 ( m + 2 ) } ( 1, - m - 1 ) C _ { - m - 1 } ^ { m + 2 } = 0 ~ }\end{array}
\end{equation}


In the metric formulation we can express the higher
spin field of arbitrary spinwith$s$


$C . { \pi } ^ { \prime }$
it is required$| n | \leq s - 1$this iteratively leads to relation
of$C . { n n } ^ { \mathrm { t r i t } + 1 }$and$C . { 0 } ^ { 1 }$ and from the$( )$analogously for $C . { - m } ^ { n + 1 }$
and $C . { 0 }$ The general form of the $C . { \pm } ^ { s }$ is then given in
terms of$C . { \pm m } ^ { m + 1 }$and coefficients$g . { \mathfrak { m } } ^ { t, \mathfrak { m } } ( \mathfrak { m }, \mathfrak { n } )$ Knowing$C . { \pm } ^ { * }$
and$C . { \pm n + 1 } ^ { r m + 1 }$allows to obtain


where$\dot { \bar { E } } . { \mu s } = \dot { \bar { A } } . { \mu } - \bar { A } . { \mu }$ and$\bar { A } . { \mu }$ and$\bar { A } . { \mu }$ we define below.
The dreibein is determined trom the background AdS
metric ()


\begin{equation}
\begin{array} { c c } { \phi _ { \mu _ { 1 },..., \mu _ { s } } = t r \left ( \vec { e } _ { ( \mu _ { 1 } }... \vec { e } _ { \mu _ { s - 1 } } \vec { E } _ { \mu _ { s } } ) \right ) ~ }\end{array}
\end{equation}


\begin{equation}
\begin{array} { r c } { ( \delta C ) _ { 0 } ^ { 1 } = \sum _ { n = 1 } ^ { s } f _ { \pm } ^ { s, n } ( \lambda ) \partial _ { z } ^ { n - 1 } \Lambda ^ { ( \ast ) } \partial _ { z } ^ { s - n } \phi \qquad } & { ( 2 ) } \\ { \dots \qquad . } \end{array}
\end{equation}


\begin{equation}
\begin{array} { l } { \epsilon _ { \bar { z } } = \frac { \varepsilon } { 2 } e ^ { \rho } ( L _ { 1 } - L _ { - 1 } ) = \frac { \varepsilon } { 2 } e ^ { \rho } ( V _ { 1 } ^ { 2 } - V _ { - 1 } ^ { 2 } ) \qquad ( 1 2 ) } \\ { \epsilon _ { \rho } = L _ { 0 } = V _ { 0 } ^ { 2 }. \qquad ( 1 3 ) } \end{array}
\end{equation}


\begin{equation}
\begin{array} { c } { e _ { z } = \frac 1 2 e ^ { p } ( L _ { 1 } + L _ { - 1 } ) = \frac 1 2 e ^ { p } ( V _ { 1 } ^ { 2 } + V _ { - 1 } ^ { 2 } ) \qquad \mathrm { ( 1 1 ) } } \\ { \qquad }\end{array}
\end{equation}


IOJ$\phi \equiv C . { 0 } ^ { 1 }$ and$f . { \pm } ^ { n, n } ( \lambda )$expressed in terms of coeffici
$\mathfrak { g } . { 1 4 } ^ { \bullet t } ( m, n )$Using tne replacement$\partial . { \rho } \rightarrow - \{ 1 \pm \lambda \}$ and writing explicitly first few n values for $f . { \pm } ^ { n, n } ( \lambda )$ allows to determine its general expression


The invariance of the equation () under the gauge transformation for${ \mathfrak { h s } } [ \lambda ] { \textcircled { \downarrow } } { \mathfrak { h s } } [ \lambda ]$for the fields A means


\begin{equation}
\begin{array} { c } { \cdot } \\ { f _ { \pm } ^ { a, n } ( \lambda ) = ( - 1 ) ^ { s } \frac { \Gamma ( s + \lambda ) } { \Gamma ( s - n + 1 \pm \lambda ) } \frac { 1 } { 2 ^ { n - 1 } ( 2 ( \frac { n } { 2 } - 1 ) )!! \left ( \frac { n - 1 } { 2 } \right )! } } \\ { \quad ~ }\end{array}
\end{equation}


Since A parameter is chiral it means$\Lambda = 0$and the field
$\bar { A }$ is essentially unchanged. The field$\bar { A } . { \mu }$is then


\begin{equation}
\begin{array} { l } { A \rightarrow A + d \Lambda + [ A, \Lambda ] _ { \star } \equiv \bar { A } \qquad \qquad ( \mathbf { 1 4 } ) } \\ { \bar { A } \rightarrow \bar { A } + d \bar { \mathbf { A } } + \left [ \bar { A }, \bar { \mathbf { \Lambda } } \right ] _ { \star } \equiv \bar { \bar { A } }. \qquad \qquad ( \mathbf { 1 5 } ) } \end{array}
\end{equation}


Substituting O in () one obtains the variation of the.
scalar field


dA reads


\begin{equation}
{ \ddot \mathcal { A } } = A _ { A d S } + d \Lambda + [ A _ { A d S }, \Lambda ] _ { * } ~. \qquad \qquad ( 1 6 )
\end{equation}


\begin{equation}
\begin{array} { r l } & { \! }\end{array}
\end{equation}


\begin{equation}
\begin{array} { r } { \mathrm { r e s - ~ o n s e r e, } } \\ { d \Lambda = \sum _ { n = 1 } ^ { 2 s - 1 } \frac { 1 } { ( n - 1 )! } V _ { s - n } ^ { s } e ^ { \left ( s - n \right ) \rho } | ( - \partial ) ^ { n - 1 } \partial \Lambda ^ { ( s ) } ( z, \bar { z } ) d z \quad ( 1 7 ) } \\ { + ( - \partial ) ^ { n - 1 } \bar { \partial } \Lambda ^ { ( s ) } ( z, \bar { z } ) d \bar { z } + \left}\end{array}\right.
\end{equation}


and


To obtain the linearised equation of motion for the scalar
field we act on ( with KG operator ().This can be
written as


To consider the coefficient in front, we focus on the term
with the lowest number of derivatives on the gauge$\dot { u } . { x }$
field$\mathbf { \Lambda } ^ { \{ \pmb { u } \} }$ obtained for n=1. Then,() becomes


\begin{equation}
\begin{array} { r } { ( \delta C ) _ { 0 } ^ { 1 } | _ { n = 1 } = ( - 1 ) ^ { s } \Lambda ^ { ( s ) } \partial ^ { s - 1 } C _ { 0 } ^ { 1 }. \qquad \qquad \qquad ( 5 ) } \\ { \mathrm { ~ h m ~ l a c u r i o n d ~ c o m t o m ~ f e r e n t i o n ~ f a r e n t i o n ~ f a r e n t i o n ~ f a r e n t i o n } } \end{array}
\end{equation}


\begin{equation}
\begin{array} { c } { \left [ A _ { A d S. } \Lambda \right ] _ { \star } = [ c ^ { \rho } V _ { 1 } ^ { 2 } d z + V _ { 0 } ^ { 2 } d \rho, ~ }\end{array}
\end{equation}


Taking$\partial . { r } \rightarrow \{ \mathbf { l } \pm \lambda \}$ in $f . { \pm } ^ { n, n } ( \lambda )$ we have taken and considering the term with hignest number of derivatives on
$C . { 0 } ^ { 1 }$ leads to


\begin{equation}
\begin{array} { c c } { \square _ { K G } \bar { C } _ { 0 } ^ { 1 } = \square _ { K G } C _ { 0 } ^ { 1 } + \square _ { K G } \delta C _ { 0 } ^ { 1 }. } & { ( 6 ) } \\ { \cdots } & { \cdots } & { \cdots } \end{array}
\end{equation}


To read out the coupling we focus on$\bar { z }, \dots \bar { z }$ component of
the field $( ^ { \cdot 1 } . { 0 }$ with lowest number of derivatives on gauge
field $\mathbf { \Lambda } ^ { ( * ) }$ The$^ \star$ multiplication of the dreibeins in. $( )$
in that case contributes only with first. $\mathfrak { g } . { u } ^ { * l } ( m, n ; \lambda )$ COefficient with the each following dreibein that is being
multiplied. More explicitly


\begin{equation}
\begin{array} { c } { \dot { \qquad }}\end{array}
\end{equation}


\begin{equation}
\begin{array} { l c } { { \square _ { K G } \big | _ { \mathrm { h i g h e s t ~ n u m b e r ~ o f ~ d e r i v a t i v e s } } ( \delta C ) _ { 0 } ^ { 1 } = } } & { { ( 7 ) } } \\ { { = ( - 1 ) ^ { s } 4 e ^ { - 2 p } \partial ( \bar { \partial } \Lambda ^ { ( s ) } \partial ^ { ( s - 1 ) } C _ { 0 } ^ { 1 } ) } } & { { ( 8 ) } } \\ { { = ( - 1 ) ^ { s } 4 e}}\end{array}
\end{equation}


on A will appear for lowest n, i.e. for$\mathfrak { n } = 1$in summation
().Knowing the relation for the trace of higher spin
generators, the required generator $V . { \mathbf { s } - \pi } ^ { * }$ will than

\end{document}
